\section{Limitations and related work}

The cross-organization run of Section~\ref{sec:eval} put the two kernels in
separate processes on one host, not on two separately administered machines.
That removes the shared address space, store, and key material, which are what
the claim is about; the shared operating system and clock source remain. The
rule does not depend on that experiment; the strength of the demonstration
does.

Section~\ref{sec:eval}'s alternative is one construction of those formats on one
corpus: what it measures is what an operator must add, not a bound on what the
parts reach.
Continuation consumption is at-most-once, so a crash between consume and
dispatch loses the call rather than repeating it: right for a refund, an
avoidable retry for an idempotent read. The chain closes in the kernel for one
hop: the next continuation in a longer chain is minted outside it.
Denial paths below the verifier still return a rendered detail alongside the
code. Receipts batch into
Merkle checkpoints in the RFC 6962 shape \cite{rfc6962} with offline inclusion
proofs, and that is the whole claim: local audit evidence with verifiable
inclusion, not a public log, and with no anti-equivocation a kernel signing two
contradictory checkpoints is detectable only by someone holding both.

\paragraph{Related work.}
Receiver-rooted policy and attenuated delegation are established by trust
management \cite{spki,keynote}, proof-carrying authorization \cite{bauer2002},
end-to-end authorization \cite{howell2000}, and macaroons \cite{macaroons}.
Modern capability formats \cite{ucan,biscuit} also support application-defined
restrictions. These systems are relevant alternatives; their different wire
shapes do not imply that they cannot express a Chio admission policy.
Authenticated Workflows \cite{authenticatedworkflows2026} also uses signed
completion attestations to enforce workflow dependencies. Chaining outcome
evidence alone is therefore not our novelty claim.

CapLease \cite{caplease2026} studies repeated agent actions under freshly issued
tokens. Its equally stateful server-ledger baseline obtains the same replay
protection under matched centralized assumptions. Durable authorization state
and an idempotent effect sink therefore belong in our comparison, not only
token identifiers. Signatures and DSSE \cite{dsse,intoto} establish authorship
and binding, while transparency systems \cite{rfc6962,scitt} address publication
and equivocation under additional assumptions.

\paragraph{The result still needed.}
Local re-execution checks a source repair without its producer's attestation,
but cannot prove private publication history. The matched ledger reproduces
these guarantees. Neither artifact validity nor signatures establish scarce
work or Sybil resistance.
Contingent payment additionally requires fair exchange~\cite{zkcp2017}; a hash
lock does not verify a hidden repair.
Useful execution across independently administered receivers, with fewer
recovery interventions or less application-specific security code than an
equally provisioned alternative, remains unmeasured. Finite checking does not
establish arbitrary correctness. Signatures do not make dishonest receivers'
external-effect claims true.
